\documentclass[11pt]{article}

% ---------- packages ----------
\usepackage[utf8]{inputenc}
\usepackage[T1]{fontenc}
\usepackage{amsmath, amssymb, amsthm}
\usepackage[dvipsnames, svgnames]{xcolor}
\usepackage{graphicx}
\usepackage{hyperref}
\usepackage{booktabs}
\usepackage{microtype}
\usepackage{geometry}
\usepackage{enumitem}
\geometry{a4paper, margin=2.5cm}
\usepackage{mathtools}
\usepackage{bm}
\usepackage{algorithm}
\usepackage{algpseudocode}
\usepackage{subcaption}
\usepackage{float}
\usepackage{framed}
\usepackage{fontawesome5}

\makeatletter
\newcommand{\github}[1]{%
   \href{#1}{\faGithubSquare}%
}
\makeatother

% ---------- hyperref config ----------
\hypersetup{
  colorlinks = true,
  linkcolor  = blue!60!black,
  citecolor  = teal,
  urlcolor   = teal,
  pdftitle   = {Arrowspace Latent Diffusion with Spectral Chart Conditioning},
  pdfauthor  = {Lorenzo Moriondo},
}

% ---------- theorem environments ----------
\newtheorem{definition}{Definition}
\newtheorem{proposition}{Proposition}
\newtheorem{remark}{Remark}
\newtheorem{conjecture}{Conjecture}

% ---------- custom commands ----------
\newcommand{\ESDM}{\textsc{esdm}}
\newcommand{\ALDSC}{\textsc{ald-sc}}
\newcommand{\VDT}{\textsc{vdt}}
\newcommand{\Lf}{L_F}
\newcommand{\Uq}{U_q}
\newcommand{\Lq}{\Lambda_q}
\newcommand{\ed}{\lambda^{\mathrm{ED}}}
\newcommand{\R}{\mathbb{R}}

% ============================================================
\begin{document}
% ============================================================

\title{%
  \textbf{Arrowspace Latent Diffusion with\\ Spectral Chart Conditioning\\
  \large \normalfont \ALDSC~A Minimal Spectral Latent Diffusion Model}
}

\author{%
  Lorenzo Moriondo \\[2pt]   \small tuned.org.uk\\   \small ORCID: 0000-0002-8804-2963\\   \texttt{tunedconsulting@gmail.com} }  \date{July 2026 \\
    \href{https://github.com/tuned-org-uk/arrowspace-latent-diffusion}{%
     \textsc{\Large\faGithub} tuned-org-uk/arrowspace-latent-diffusion} \\
     \href{https://github.com/tuned-org-uk/entropic-semantic-diffusion}{%
     \textsc{\Large\faGithub} tuned-org-uk/entropic-semantic-diffusion}
  }
    \maketitle

    \begin{abstract}
    We present \ALDSC~(\emph{Arrowspace Latent Diffusion with Spectral
    Chart Conditioning}), a spectral latent diffusion model in which
    decoding is performed on the feature-space manifold defined by a
    frozen ArrowSpace graph Laplacian $L_F$ and its associated
    energy-dispersion network $\lambda^{\mathrm{ED}}$.

    The model encodes each image into a 2.5-D representation: a spatial
    VAE latent $z$ for local detail, and a feature field $A$ projected
    onto the smooth eigenspace of $L_F$. The 2.5-D encoding target is
    the DualSpaceMatrix $M_N = \alpha\|VV^\top\|_F - \beta\|V L_F
    V^\top\|_F$, which fuses item-space geometry with feature-space
    topology. From the projected feature field we extract a compact
    spectral chart $c_{\mathrm{spec}} \in \mathbb{R}^{3q}$ comprising
    normalised band energies, the projected dispersion prior, and the
    Laplacian eigenvalues.

    Diffusion operates only on $z$ with v-prediction and a cosine
    schedule; $c_{\mathrm{spec}}$ is a conditioning side-channel, not a
    second diffused state. The DiT denoiser receives $c_{\mathrm{spec}}$
    via adaptive layer normalisation, with classifier-free guidance
    dropout.

    The decoder is the research contribution: a \emph{graph-structured}
    decoder in which $L_F$ (via its eigenvectors $U_q$) defines the
    reconstruction paths and $\lambda^{\mathrm{ED}}$ (via
    $c_{\mathrm{spec}}$) gates the energy allocation at each resolution.
    The \emph{WaveReconstructionBlock} projects feature activations onto
    the chart basis, gates them by the dispersion-derived weights, and
    lifts them back --- the graph-theoretic analogue of the VAE
    reparameterization trick, replacing unconstrained convolutions with
    propagation along the graph's smooth directions.

    The Barontini entropic clock is an active part of both inference and
    decoding. The per-mode schedule $\tau_k(t) = \nu_k \cdot t$ and the
    heat-death criterion $\sum_k \nu_k \bar\alpha_k(t) < \varepsilon$
    terminate sampling intrinsically rather than at a fixed step count.
    The \emph{ClockGatedGraphDecoder} modulates decoding tempo by
    $\bar\alpha_k(t)$: graph-structured reconstruction effort increases
    as the denoising process resolves each spectral mode.

    The full pipeline --- prior construction, spectral VAE, graph
    decoder, v-prediction diffusion training, DDIM/Euler sampling with
    spectral stopping, and end-to-end image generation --- is
    implemented in PyTorch and validated by 107 unit tests.

    The central falsifiable claim is that decoding on the feature-space
    manifold yields better global semantic coherence under compression
    than decoding on an unconstrained ambient latent.
    \end{abstract}

    \vspace{1em} \noindent \textbf{Keywords:} latent diffusion, ArrowSpace, graph Laplacian, spectral conditioning, feature-space geometry, v-prediction, classifier-free guidance, thermal time.  \newpage \tableofcontents \newpage

    % ============================================================
    \section{Introduction} \label{sec:introduction}
    % ============================================================
    Latent diffusion models achieve high-quality image synthesis by
    moving the iterative denoising process from pixel space to a
    compressed latent space~\cite{Rombach2022}. While effective, this
    design leaves the semantic structure of the latent space largely
    unconstrained: smoothness, coherence, and global feature
    organisation emerge only indirectly from reconstruction and
    denoising losses.

    The predecessor \ESDM~\cite{Moriondo2026ESDM_code} proposed a
    spectrally conditioned diffusion architecture in which a frozen
    ArrowSpace graph Laplacian~\cite{Moriondo2025JOSS,
    MoriondoAzizi2026} defines the valid semantic subspace, while a
    full entropic clock --- wave recurrence, vibrational pump, and a
    bright--dark KL entropy clock inspired by Barontini's cold-atom
    experiment~\cite{Barontini2026} --- drives the diffusion schedule
    from internal thermodynamics. That model is theoretically rich
    but implementation-heavy.

    This paper strips \ESDM~down to its testable core. We retain the
    frozen ArrowSpace prior and the 2.5-D latent (a spatial latent
    plus a spectral chart) but replace the entropic clock with a
    standard cosine schedule and v-prediction. The spectral chart
    enters as a \emph{conditioning side-channel}, not as a second
    diffused state. The result is a model that is recognisably a
    latent diffusion system but whose semantic organisation is
    explicitly governed by graph Laplacian structure, trainable on a
    single CPU in minutes, and implementable in under 700 lines.
    
        \begin{framed}
        \noindent\textbf{What is implemented.} The code at
        \texttt{arrowspace-latent-diffusion} implements:
        \begin{itemize}[leftmargin=1.5em]
        \item The frozen ArrowSpace prior $(\Lf, \Uq, \Lq, \ed)$ via the
              ArrowSpace library or a kNN fallback (\texttt{wire\_graph.py},
              \texttt{arrow\_prior.py}, \texttt{build\_prior.py}).
        \item The 2.5-D encoding target: the DualSpaceMatrix
              $M_N = \alpha\|VV^\top\|_F - \beta\|V \Lf V^\top\|_F$
              (\texttt{dual\_space.py}).
        \item The spectral VAE with dual-head encoder producing $(z, A)$
              and $c_{\mathrm{spec}} \in \R^{3q}$
              (\texttt{vae.py}, \texttt{losses.py}).
        \item The DiT denoiser with AdaLN conditioning on $c_{\mathrm{spec}}$
              and classifier-free guidance dropout (\texttt{dit.py}).
        \item Cosine and linear noise schedules with v-prediction
              (\texttt{schedule.py}).
        \item A graph-structured decoder where $L_F$ (via $\Uq$) defines
              reconstruction paths and $\ed$ (via $c_{\mathrm{spec}}$) gates
              energy allocation. The \texttt{WaveReconstructionBlock}
              propagates information along the graph's smooth directions ---
              the graph-theoretic analogue of the VAE reparameterization trick
              (\texttt{graph\_decoder.py}).
        \item The Barontini entropic clock as an active intrinsic stopping
              criterion: the sampler terminates when
              $\sum_k \nu_k \bar\alpha_k(t) < \varepsilon$
              (\texttt{spectral\_schedule.py}, \texttt{sampling.py}).
              The \texttt{ClockGatedGraphDecoder} also modulates decoding
              tempo by $\bar\alpha_k(t)$.
        \item DDIM and Euler samplers, end-to-end image generation via CLI
              (\texttt{sampling.py}, \texttt{scripts/sample.py}).
        \end{itemize}
        The full entropic-clock formulation --- wave recurrence,
        vibrational pump, density matrices, and the bright--dark KL clock
        --- is implemented in the predecessor repository
        \texttt{entropic-semantic-diffusion} and is cited here as
        background. The code is validated by 107 unit tests on CPU.
        \end{framed}

    The three commitments of \ALDSC~are:
    \begin{enumerate}[leftmargin=1.5em]
      \item a frozen ArrowSpace prior $(\Lf,\Uq,\Lq,\ed)$ built once from corpus feature statistics;
      \item a 2.5-D latent consisting of a standard spatial latent plus a compact spectral chart $c_{\mathrm{spec}}$;
      \item a decoder and denoiser conditioned by feature-space spectral geometry.
    \end{enumerate}

    % ============================================================
    \section{Background}
    \label{sec:background}
    % ============================================================

    \subsection{ArrowSpace and feature-space geometry}
    \label{sec:arrowspace}
    ArrowSpace~\cite{Moriondo2025JOSS} introduces a graph-based
    treatment of vector spaces in which feature dimensions, rather
    than only data points, participate in a structural semantic
    geometry. Given a corpus of embeddings with feature dimension $F$,
    one constructs a feature-space graph and its Laplacian $\Lf \in
    \R^{F \times F}$. The eigenvectors of $\Lf$ define smooth and
    rough directions in feature space, while the Energy Dispersion
    Networks formulation~\cite{MoriondoAzizi2026} treats the feature
    columns of an embedding matrix as nodes in a semantic energy
    network.

    In \ALDSC~we treat ArrowSpace as a frozen semantic prior:
    \[ \mathcal{G}_{\mathrm{Arrow}} = (\Lf,\Uq,\Lq,\ed),\]
    where $\Uq \in \R^{F \times q}$ contains the retained smooth
    modes, $\Lq = \mathrm{diag}(\nu_1,\ldots,\nu_q)$ are the
    corresponding eigenvalues, and $\ed \in [0,1]^F$ is an empirical
    energy-dispersion distribution over feature nodes. This prior is
    built once from the training corpus and never updated during
    training (all components are stored as non-trainable buffers).

    \subsection{Latent diffusion}
    \label{sec:latent_diffusion}
    Latent diffusion models~\cite{Rombach2022} shift the forward and
    reverse diffusion processes from pixel space into a learned
    compressed latent representation. An image $x \in \R^{H \times W
    \times 3}$ is encoded into a latent tensor $z \in \R^{h \times w
    \times c}$, denoised over a sequence of timesteps, and decoded
    back to RGB. Standard latent diffusion assumes that the VAE latent
    is sufficient as the sole state on which the generative process
    must operate. Our claim is that this is lossy: a spatial latent is
    effective for local visual detail, but a separate feature-space
    spectral state is needed to preserve global semantic organisation
    under compression~\cite{MoriondoAzizi2026}.

    \subsection{Internal time and entropic motivation}
    \label{sec:internal_time}
    The original \ESDM~framing was motivated by the problem of time
    in closed systems. In~\cite{ConnesRovelli1994} it is argued that
    in a generally covariant theory, the physical time-flow may arise
    from the thermodynamical state of the system rather than from a
    universal external parameter. Barontini's cold-atom
    experiment~\cite{Barontini2026} provides an operational testbed:
    a closed system partitioned into observed and unobserved sectors
    admits an entropic time constructed from internal entropy
    exchange. We retain this as a design principle --- generative
    dynamics should be informed by state-derived geometry --- but in
    \ALDSC~it is instantiated through a frozen spectral prior rather
    than through a fully entropic recurrence. The full entropic clock
    is deferred to Section~\ref{sec:entropic_clock} as a theoretical
    extension.

    % ============================================================
    \section{Design transition: from \ESDM~to \ALDSC}
    \label{sec:transition}
    % ============================================================

    Earlier \ESDM~drafts coupled five ingredients tightly: a
    dual-space semantic operator (ArrowSpace), a wave recurrence over
    item states, a bright--dark sector entropy clock, an entropic
    pump, and a topology decoder acting as an inverse spectral map.
    This unified picture was conceptually rich but
    implementation-heavy.

    \begin{remark}
    The strongest reusable result of \ESDM~is not the full
    second-order entropic dynamics but the geometric constraint that
    generation should be restricted to a smooth ArrowSpace subspace
    and decoded relative to a feature-space manifold.
    \end{remark}

    \ALDSC~therefore keeps:
    \begin{itemize}[leftmargin=1.5em]
    \item the frozen ArrowSpace prior $(\Lf,\Uq,\Lq,\ed)$;
    \item the projection onto smooth spectral modes $\Pi_q = \Uq\Uq^\top$;
    \item the 2.5-D representation (spatial latent + spectral chart);
    \item the heat-kernel interpretation of the spectral chart;
    \item topology-adaptive decoding via spectral gates.
    \end{itemize}

    It \textbf{defers} to the predecessor \ESDM~codebase:
    \begin{itemize}[leftmargin=1.5em]
    \item the full wave recurrence;
    \item explicit Rayleigh-gradient restoring forces;
    \item the bright--dark KL clock as the primary training schedule;
    \item the entropic pump as a learned forcing term.
    \end{itemize}

    The entropic clock is implemented as an intrinsic stopping criterion
    in Section~\ref{sec:entropic_clock}: the sampler terminates when
    the spectrally-weighted heat-death metric drops below a threshold.

    % ============================================================
    \section{The 2.5-D latent}
    \label{sec:25d}
    % ============================================================

    \subsection{Dual representation}
    \label{sec:dual_representation}
    Let $x \in \R^{H \times W \times 3}$ be an image. The encoder
    produces two coupled outputs:
    \[x \xrightarrow{E_\phi} (z, A),\]
    where $z \in \R^{h \times w \times c}$ is the spatial latent and
    $A \in \R^{N \times F}$ is a feature field ($N = h \cdot w$).

    \noindent\textbf{Implementation:} \texttt{vae.py} --
    \texttt{SpectralVAE.encode()} produces $(z, A, \mu, \mathrm{logvar})$ via
    a shared conv trunk, a spatial head (Conv2d $\to$ $\mu$, logvar), and a
    feature head (adaptive pool $\to$ linear projection to $F$ dimensions).

    \subsection{Smooth projection}
    \label{sec:smooth_projection}
    The ArrowSpace prior defines a smooth subspace through the
    projection $\Pi_q = \Uq \Uq^\top$. Applying this to the feature
    field yields $A_{\parallel} = A \Pi_q$ and $A_{\perp} = A(I -
    \Pi_q)$. The parallel component lies in the smooth semantic
    subspace; the perpendicular captures off-manifold residual.

    \noindent\textbf{Implementation:} \texttt{arrow\_prior.py} --
    \texttt{ArrowSpacePrior.project\_to\_chart(A)} computes $A \Uq
    \Uq^\top$. The prior stores $L_F$, $U_q$, eigenvalues, and
    energy-dispersion as non-trainable buffers (zero
    \texttt{nn.Parameter}).

    \subsection{Spectral chart}
    \label{sec:spectral_chart}
    From the projected feature field we extract chart coordinates
    $s = \mathrm{Pool}(A_{\parallel}) \Uq \in \R^q$, per-mode band
    energies $e_k = \frac{1}{N}\|A u_k\|_2^2$, normalised energies
    $\tilde e_k = e_k / (\sum_j e_j + \varepsilon)$, the projected
    energy-dispersion prior $\ed_{k^{\mathrm{chart}}} = \sum_f \ed_f
    U_{f,k}^2$, and the full conditioning vector
    \[c_{\mathrm{spec}} = [\tilde e_1,\ldots,\tilde e_q,\;
    \ed_{1^{\mathrm{chart}}},\ldots,\ed_{q^{\mathrm{chart}}},\;
    \nu_1,\ldots,\nu_q] \in \R^{3q}.\]

    \noindent\textbf{Implementation:} \texttt{arrow\_prior.py} --
    \texttt{chart\_energy\_descriptor(A)} returns $c_{\mathrm{spec}}$
    of shape $(B, 3q)$. No MLP is applied in the current
    implementation; the raw $3q$ vector is used directly as
    conditioning. An MLP can be added without changing the interface.

    % ============================================================
    \section{Model architecture}
    \label{sec:model}
    % ============================================================

    \subsection{Overview}
    \label{sec:overview}
    \ALDSC~consists of four components:
    \begin{enumerate}[leftmargin=1.5em]
    \item a frozen ArrowSpace prior $(\Lf,\Uq,\Lq,\ed)$;
    \item a spectral VAE encoder producing $z$ and $A$;
    \item a latent diffusion denoiser (DiT) operating only on $z$;
    \item a topology-adaptive decoder conditioned on $c_{\mathrm{spec}}$.
    \end{enumerate}
    The full pathway is
    \[x \;\xrightarrow{E_\phi}\; (z, A) \;\xrightarrow{\text{chart}}
    (z, c_{\mathrm{spec}}) \;\xrightarrow{\text{diffusion}} \hat z
    \;\xrightarrow{D_\psi(\cdot \mid c_{\mathrm{spec}})} \hat x.\]

    \subsection{Encoder}
    \label{sec:encoder}
    The encoder is a shared visual trunk with two heads. The spatial
    head produces $z \sim q_\phi(z \mid x)$ via mean and log-variance
    fields. The feature head produces $A = E_A(x) \in \R^{N \times F}$.
    The feature field is not diffused; it is projected onto the
    ArrowSpace subspace and summarised into $c_{\mathrm{spec}}$.

    \subsection{Diffusion backbone}
    \label{sec:backbone}
    We apply diffusion only to the spatial latent:
    \[q(z_t \mid z_0) = \mathcal{N}(\sqrt{\bar\alpha_t} z_0,\;
    (1-\bar\alpha_t)I).\]
    A DiT denoiser predicts velocity:
    \[v_\theta(z_t, t, c_{\mathrm{text}}, c_{\mathrm{spec}}).\]
    The spectral chart enters via adaptive layer normalisation
    (AdaLN): the conditioning vector $c_{\mathrm{spec}}$ is fused with
    the timestep embedding and projected to per-block scale and shift
    parameters. Classifier-free guidance (CFG) dropout on
    $c_{\mathrm{spec}}$ is applied during training.

    \noindent\textbf{Implementation:} \texttt{dit.py} --
    \texttt{MinimalDiT} with patchify, sinusoidal time embedding, AdaLN
    blocks, and CFG dropout. \texttt{schedule.py} provides
    \texttt{CosineSchedule} and \texttt{LinearSchedule} with
    \texttt{add\_noise}, \texttt{v\_target}, \texttt{sample\_batch}.

    \subsection{Graph-structured decoder}
    \label{sec:decoder}
    The decoder reconstructs pixels by propagating information along the
    ArrowSpace graph's smooth directions, gated by the dispersion
    network. This replaces the standard VAE decoder's unconstrained
    convolutions with a constructive decoding operator where $L_F$ (via
    $U_q$) defines the reconstruction paths and $\lambda^{\mathrm{ED}}$
    (via $c_{\mathrm{spec}}$) defines the energy allocation.

    At each resolution stage, the \emph{WaveReconstructionBlock}
    performs:
    \begin{enumerate}
    \item Pool feature activations: $A = \mathrm{Pool}(H^{(\ell)})$
    \item Project to chart: $\hat H = A \Uq$ (decode along smooth directions)
    \item Gate by dispersion: $g = \sigma(W \, c_{\mathrm{spec}})$
    \item Lift back: $A' = (\hat H \odot g) \Uq^\top$ (reconstruct in feature space)
    \item Residual update: $H^{(\ell+1)} = H^{(\ell)} + \mathrm{Conv}(\mathrm{norm}(H^{(\ell)} + \Delta))$
    \end{enumerate}
    where $\Delta$ is the feature-space reconstruction projected back to
    channel dimension. This is the graph-theoretic analogue of the VAE
    reparameterization trick: information flows along the graph's
    eigenvectors rather than through unconstrained learned weights.

    \noindent\textbf{Clock-gated tempo.} The \emph{ClockGatedGraphDecoder}
    modulates the strength of each wave block by the spectral schedule's
    $\bar\alpha_k(t)$: early in the denoising process (high noise), gates
    are weak (minimal graph-structured reconstruction); late (low noise),
    gates are strong (full graph-structured reconstruction). This
    implements the Barontini principle within the decoder: the graph
    geometry governs \emph{when} reconstruction effort is allocated.

    \noindent\textbf{Implementation:} \texttt{graph\_decoder.py} --
    \texttt{GraphDecoder} uses 3 \texttt{WaveReconstructionBlock}
    instances at decreasing channel resolution ($4\mathrm{ch}
    \to 2\mathrm{ch} \to \mathrm{ch}$).
    \texttt{ClockGatedGraphDecoder} adds the
    \texttt{SpectralSchedule} for tempo modulation. The
    \texttt{WaveReconstructionBlock} stores $U_q$ as a frozen buffer
    (zero trainable parameters in the graph structure itself). The
    simplified single-gate decoder in \texttt{vae.py} remains for
    backward compatibility.

    % ============================================================
    \section{Heat-kernel prior}
    \label{sec:heat_kernel}
    % ============================================================
    The spectral chart admits a natural diffusion-geometric prior. Let
    $s = (s_1,\ldots,s_q)^\top$ denote the chart coordinates. We
    define
    \[p_t(s) = \mathcal{N}\!\left( 0,\;
    \mathrm{diag}\bigl[w_k e^{-2t\nu_k}\bigr]_{k=1}^q \right),\]
    where $w_k = \frac{\exp(-\rho \ed_{k^{\mathrm{chart}}})} {\sum_j
    \exp(-\rho \ed_{j^{\mathrm{chart}}})}$. Low eigenvalue modes
    persist under diffusion; high eigenvalue modes decay rapidly. The
    energy-dispersion weights bias the chart toward regions of
    spectral space occupied by corpus data.

    \begin{definition}[2.5-D manifold prior]
    The pair $(z,s)$ defines a 2.5-D manifold prior if $z$ carries
    local spatial degrees of freedom and $s$ is drawn from a
    heat-kernel distribution on the ArrowSpace spectral chart.
    \end{definition}

    % ============================================================
    \section{Objective}
    \label{sec:objective}
    % ============================================================
    The training objective is
    \[\mathcal{L} = \mathcal{L}_{\mathrm{diff}} +
    \lambda_{\mathrm{rec}}\mathcal{L}_{\mathrm{rec}} +
    \lambda_{\mathrm{chart}}\mathcal{L}_{\mathrm{chart}} +
    \lambda_{\mathrm{smooth}}\mathcal{L}_{\mathrm{smooth}} +
    \lambda_{\mathrm{kl}}\mathcal{L}_{\mathrm{kl}}.\]

    \subsection{Latent diffusion loss}
    Using velocity prediction,
    \[\mathcal{L}_{\mathrm{diff}} = \mathbb{E}_{z_0,\epsilon,t}
    \left[ \left\| v - v_\theta(z_t,t,c_{\mathrm{spec}}) \right\|_2^2
    \right].\]

    \subsection{Reconstruction and VAE regularisation}
    \[\mathcal{L}_{\mathrm{rec}} = \|x-\hat x\|_1 +
    \lambda_{\mathrm{perc}}\mathcal{L}_{\mathrm{perc}}(x,\hat x),\qquad
    \mathcal{L}_{\mathrm{kl}} = D_{\mathrm{KL}}\!\bigl(q_\phi(z \mid
    x)\,\|\,p(z)\bigr).\]

    \subsection{Spectral consistency}
    Re-encoding $\hat x$ produces $\hat A$; we require the spectral
    band allocation to match:
    \[\mathcal{L}_{\mathrm{chart}} = \left\| \tilde e(A) - \tilde
    e(\hat A) \right\|_2^2.\]

    \subsection{Off-manifold penalty}
    \[\mathcal{L}_{\mathrm{smooth}} = \frac{
    \|A(I-\Uq\Uq^\top)\|_F^2 }{ \|A\|_F^2 + \varepsilon }.\]

    \noindent\textbf{Implementation:} \texttt{losses.py} --
    \texttt{ALDSCLoss} computes all five terms. The VAE KL is
    included in the current implementation (\texttt{vae.kl\_loss()}
    extracts $\mu$ and logvar from the last forward pass).

    % ============================================================
    \section{Entropic time from the ArrowSpace spectrum}
    \label{sec:entropic_clock}
    % ============================================================

    \noindent\textbf{Implementation note.} The Barontini-inspired
    entropic clock is an active part of the \ALDSC~inference pipeline.
    The heat-death stopping criterion (Section~\ref{sec:heat_death_clock})
    is wired into the DDIM and Euler samplers, terminating sampling
    when $\sum_k \nu_k \cdot \bar\alpha_k(t) < \varepsilon$.

    The heat-kernel prior admits a natural reading in terms of
    entropic time. The ArrowSpace Laplacian eigenvalues play the role
    of Barontini's barrier-controlled entropy exchange rates, yielding
    a per-mode entropic clock that requires no additional learned
    parameters and no auxiliary model. This connects the frozen
    spectral prior directly to information-theoretic advances in
    diffusion scheduling~\cite{Stancevic2025}.

    % ------------------------------------------------------------
    \subsection{Entropy rate of the heat-kernel chart}
    \label{sec:entropy_rate}
    % ------------------------------------------------------------
    The per-mode variance under the heat-kernel prior is
    $\sigma_k^2(t) = w_k e^{-2t\nu_k}$. The Shannon entropy of a
    Gaussian with this variance is $S_k(t) = \tfrac{1}{2}\log(2\pi e\,
    w_k) - t\nu_k$, so the entropy rate is
    \[\boxed{\frac{dS_k}{dt} = -\nu_k.}\]
    Each Laplacian eigenvalue $\nu_k$ \emph{is} the entropy exchange
    rate for mode~$k$.

    % ------------------------------------------------------------
    \subsection{Correspondence to Barontini's two-sector clock}
    \label{sec:barontini_mapping}
    % ------------------------------------------------------------
    Barontini partitions a closed Bose--Einstein condensate into a
    bright (observed) and dark (unobserved) sector, separated by a
    tunable potential barrier of height~$V$. Entropic time is
    constructed from the absolute entropy exchange between sectors.
    Table~\ref{tab:barontini} maps this onto the \ALDSC~spectral
    chart.

    \begin{table}[ht]
    \centering
    \small
    \caption{Correspondence between Barontini's BEC experiment and the
    ALD-SC spectral chart.}
    \label{tab:barontini}
    \begin{tabular}{ll}
    \toprule
    \textbf{Barontini experiment} & \textbf{ALD-SC spectral chart} \\
    \midrule
    Bright sector (observed) & Active modes: low $\nu_k$, high variance \\
    Dark sector (unobserved) & Inactive modes: high $\nu_k$, collapsed variance \\
    Barrier height $V$ & Laplacian eigenvalue $\nu_k$ (per-mode barrier) \\
    Entropy exchange rate & $|dS_k/dt| = \nu_k$ \\
    Entropic time $\tau$ & Per-mode $\tau_k(t) = \nu_k \cdot t$ \\
    Big bang / big crunch & Mode activation / collapse during denoising \\
    Heat death & All modes at equilibrium variance \\
    Barrier-height sweep & Varying $q$ (number of retained modes) \\
    \bottomrule
    \end{tabular}
    \end{table}

    % ------------------------------------------------------------
    \subsection{Per-mode entropic schedule and stopping criterion}
    \label{sec:per_mode_schedule}
    % ------------------------------------------------------------
    Instead of a single global schedule $\bar\alpha(t)$, each mode
    receives its own entropic schedule driven by its internal clock:
    \[\tau_k(t) = \nu_k \cdot t, \qquad
    \bar\alpha_k(\tau_k) = \cos^2\!\left(\frac{\pi}{2}\cdot\frac{\tau_k}{\tau_{k,\max}}\right),\]
    where $\tau_{k,\max} = \nu_k \cdot T$ for a reference horizon~$T$.

    \begin{remark}
    \textbf{Mathematical observation.} Since $\nu_k$ cancels in the
    ratio $\tau_k / \tau_{k,\max} = t/T$, all modes have identical
    $\bar\alpha_k$ in external time. The per-mode differentiation
    arises through the \emph{heat-death metric} $\sum_k \nu_k \cdot
    \mathrm{mmse}_k(t)$, which is weighted by $\nu_k$: high-$\nu_k$
    modes dominate the stopping criterion. All modes formally reach
    equilibrium at the same external time $T$, but high-$\nu_k$ modes
    accumulate entropic time $\tau_k$ faster.
    \end{remark}

    \noindent\textbf{Implementation:} \texttt{spectral\_schedule.py} --
    \texttt{SpectralSchedule} implements $\tau_k(t)$, $\bar\alpha_k(t)$,
    the entropy rate $dS_k/dt = -\nu_k$, and the heat-death stopping
    criterion $\sum_k \nu_k \cdot \bar\alpha_k(t) < \varepsilon$.
    The samplers in \texttt{sampling.py} accept an optional
    \texttt{spectral\_schedule} parameter; when provided, sampling
    stops early at heat death instead of at a fixed step count. All
    components are frozen (zero \texttt{nn.Parameter}). 21 unit tests
    verify shapes, monotonicity, boundary cases, the $\nu_k$
    cancellation, and the spectral stopping criterion.

    % ------------------------------------------------------------
    \subsection{Connection to information-theoretic schedulers}
    \label{sec:stancevic}
    % ------------------------------------------------------------
    Stancevic et al.~\cite{Stancevic2025} prove that the conditional
    entropy $\mathbf{H}[\mathbf{x}_0 \mid \mathbf{x}_t]$ is a
    schedule-invariant time reparameterisation for diffusion models.
    The entropy rate decomposes over modes when the chart covariance
    is diagonal:
    \[\dot{\mathbf{H}}[s_0 \mid s_t] = \sum_{k=1}^{q} \nu_k \cdot
    \mathrm{mmse}_k(t).\]
    The \ALDSC~contribution is that the ArrowSpace spectrum provides
    the decomposition basis a priori: the $\nu_k$ are frozen graph
    eigenvalues.

    Table~\ref{tab:scheduler_comparison} contrasts \ALDSC~with prior
    approaches.

    \begin{table}[ht]
    \centering
    \small
    \caption{Comparison of noise-scheduling approaches.}
    \label{tab:scheduler_comparison}
    \begin{tabular}{llll}
    \toprule
    \textbf{Feature} & \textbf{Stancevic} & \textbf{Spectrally-Guided} & \textbf{ALD-SC} \\
     & \textbf{et al.} & \textbf{(Esteves)} & \textbf{(this work)} \\
    \midrule
    Time variable & Cond.\ entropy & RAPSD profile & Per-mode $\tau_k$ \\
    Spectral basis & Fourier (pixel) & Fourier (pixel) & ArrowSpace (semantic) \\
    Per-mode schedule & Spectral variant & No & $\tau_k$ (stopping only) \\
    Geometry source & Image frequency & Image spectrum & Feature-space graph \\
    Online adaptation & No (post-train) & No & No (static $\nu_k$) \\
    Barrier analogue & None & None & $q$ (retained modes) \\
    \bottomrule
    \end{tabular}
    \end{table}

    % ------------------------------------------------------------
    \subsection{Heat death as a stopping criterion}
    \label{sec:heat_death_clock}
    % ------------------------------------------------------------
    In \ALDSC, heat death corresponds to all modes reaching
    equilibrium variance:
    \[\mathrm{heat\ death} \iff \sum_{k=1}^{q} \nu_k \cdot
    \mathrm{mmse}_k(t) < \varepsilon.\]
    This provides an intrinsic, data-dependent stopping criterion for
    sampling.

    % ============================================================
    \section{Minimal viable experiment}
    \label{sec:experiments}
    % ============================================================

    We implement the full \ALDSC~pipeline and validate it on a
    controlled toy experiment. The goal is not to claim
    state-of-the-art image quality but to demonstrate that the
    spectral chart pipeline is functional and that
    $c_{\mathrm{spec}}$ measurably affects generation.

    \subsection{Setup}
    \label{sec:setup}
    \begin{itemize}[leftmargin=1.5em]
    \item \textbf{Prior:} $F=32$, $q=8$, $k=4$ (kNN over feature
        columns). Built from 64 synthetic embeddings.
    \item \textbf{VAE:} 3-channel input, 4-channel latent
        ($32\times32 \to 8\times8$ latent), 32 base channels.
    \item \textbf{DiT:} patch size 2, dim 64, depth 4, 4 heads,
        $c_{\mathrm{spec}}$ dim $= 3q = 24$, CFG dropout 0.1.
    \item \textbf{Schedule:} Cosine, 1000 steps, v-prediction.
    \item \textbf{Training:} 20 epochs VAE, 20 epochs diffusion,
        Adam lr$=10^{-3}$, batch size 8, seed 3407.
    \item \textbf{Data:} 32 synthetic $32\times32$ noise images
        (ToyImageDataset).
    \end{itemize}

    \subsection{Results}
    \label{sec:results}
    \begin{itemize}[leftmargin=1.5em]
    \item \textbf{Image generation:} \texttt{scripts/sample.py}
        generates and saves a PNG end-to-end.
    \item \textbf{Conditioning sensitivity:} swapping
        $c_{\mathrm{spec}}$ between images produces measurably
        different DiT velocity predictions (verified by unit test
        after perturbing zero-initialised AdaLN weights).
    \item \textbf{CFG dropout:} at dropout probability 1.0, the model
        falls back to the unconditional path; at 0.0, $c_{\mathrm{spec}}$
        is always used.
    \item \textbf{Frozen prior:} verified that the prior has zero
        trainable parameters and that VAE parameters stay frozen
        during diffusion training.
    \item \textbf{Graph-structured decoding:} the \texttt{GraphDecoder}
        uses \texttt{WaveReconstructionBlock} instances that propagate
        information along $U_q$ directions, gated by $c_{\mathrm{spec}}$.
        The \texttt{ClockGatedGraphDecoder} modulates gate strength by
        the spectral schedule (verified: different diffusion times
        produce different outputs).
    \item \textbf{Spectral stopping:} the per-mode $\tau_k(t)$,
        $\bar\alpha_k$, entropy rate, and heat-death metric are
        computed on frozen $\nu_k$. The DDIM sampler terminates early
        when the heat-death criterion is met (verified by unit test:
        high-$\varepsilon$ stops before max steps; same seed is
        deterministic).
    \end{itemize}

    \subsection{Limitations}
    \label{sec:limitations}
    \begin{itemize}[leftmargin=1.5em]
    \item Experiments use synthetic noise data, not real images. Real
        data experiments (CIFAR-10 with DINO/SigLIP embeddings) are
        the next step.
    \item The graph decoder uses a single graph-filter step per block,
        not the full second-order wave recurrence from \ESDM.
    \item The entropic clock provides an intrinsic stopping criterion
        but not an active training schedule; the cosine schedule is
        still used during training.
    \item No FID or quantitative image-quality metrics are reported.
    \end{itemize}

    \subsection{Metrics (for future real-data experiments)}
    \label{sec:metrics}
    Beyond standard image-quality metrics (PSNR, SSIM, LPIPS, FID),
    evaluation should include spectral diagnostics:
    \begin{itemize}[leftmargin=1.5em]
    \item chart-energy error: $\|\tilde e(x) - \tilde e(\hat x)\|^2$;
    \item off-manifold energy ratio;
    \item chart interpolation smoothness;
    \item global semantic coherence under low latent capacity.
    \end{itemize}

    % ============================================================
    \section{Connection to internal time}
    \label{sec:internal_time_connection}
    % ============================================================
    Although \ALDSC~does not use the full entropic clock as its
    primary training schedule, the original motivation remains
    important. The conceptual bridge is that geometry and dynamics
    should be derived from internal state rather than imposed entirely
    from outside. In \ESDM~this took the form of an entropy-derived
    temporal variable. In \ALDSC~it takes the form of a feature-space
    geometric prior that governs which semantic directions are valid
    and how their energy is distributed. The spectral schedule
    (Section~\ref{sec:per_mode_schedule}) provides the active
    stopping criterion, making the Barontini principle operational
    within the minimal architecture.

    % ============================================================
    \section{Conclusion}
    \label{sec:conclusion}
    % ============================================================
    We presented \ALDSC, a minimal spectrally conditioned latent
    diffusion architecture grounded in ArrowSpace feature geometry.
    The model combines a standard spatial latent with a frozen
    spectral prior and a decoder conditioned on a compact ArrowSpace
    chart. This realises the central 2.5-D claim: local image detail
    lives in a spatial latent plane, while global semantic coherence
    is governed by a lower-dimensional graph manifold.

    The full \ALDSC~pipeline --- prior construction, spectral VAE,
    v-prediction diffusion, and DDIM/Euler sampling --- is implemented
    and validated by 67 unit tests. The entropic clock is included as
    an active intrinsic stopping criterion in the inference pipeline,
    making the Barontini principle operational within the minimal
    architecture. The upgrade path to the full \ESDM~remains open
    through the entropic training schedule.

    The original entropic formulation remains an important theoretical
    horizon. The present architecture offers a cleaner and more
    tractable route to testing the key scientific idea: that generative
    models should reconstruct and sample not only in compressed latent
    coordinates but on a semantic manifold defined by the graph
    Laplacian structure of the feature space itself.

    % ============================================================
    \appendix
    \section{Reproducibility: code-to-paper mapping}
    \label{sec:reproducibility}
    % ============================================================

    \begin{table}[ht]
    \centering
    \small
    \caption{Mapping from paper sections to implementation files.}
    \label{tab:reproducibility}
    \begin{tabular}{lll}
    \toprule
    \textbf{Paper section} & \textbf{File} & \textbf{Status} \\
    \midrule
    \S\ref{sec:25d}: 2.5-D latent & \texttt{vae.py}, \texttt{arrow\_prior.py} & Implemented \\
    \S\ref{sec:smooth_projection}: projection & \texttt{arrow\_prior.py} & Implemented \\
    \S\ref{sec:spectral_chart}: $c_{\mathrm{spec}}$ & \texttt{arrow\_prior.py} & Implemented \\
    \S\ref{sec:backbone}: DiT + AdaLN & \texttt{dit.py} & Implemented \\
    \S\ref{sec:backbone}: schedules & \texttt{schedule.py} & Implemented \\
    \S\ref{sec:decoder}: graph decoder & \texttt{graph\_decoder.py} & Implemented \\
    \S\ref{sec:objective}: losses & \texttt{losses.py} & Implemented \\
    \S\ref{sec:heat_kernel}: heat-kernel prior & --- & Theoretical \\
    \S\ref{sec:entropic_clock}: entropic clock & \texttt{spectral\_schedule.py}, \texttt{sampling.py} & Implemented \\
    \S\ref{sec:per_mode_schedule}: per-mode schedule & \texttt{spectral\_schedule.py} & Implemented \\
    \S\ref{sec:heat_death_clock}: heat death & \texttt{spectral\_schedule.py}, \texttt{sampling.py} & Implemented \\
    \S\ref{sec:experiments}: image generation & \texttt{scripts/sample.py} & Implemented \\
    \bottomrule
    \end{tabular}
    \end{table}

    All code is at
    \url{https://github.com/tuned-org-uk/arrowspace-latent-diffusion}.
    Run tests: \texttt{uv run pytest tests/ -v} (67 tests, CPU).
    Generate an image: \texttt{uv run python scripts/sample.py}.

    % ============================================================
    \newpage
    \begin{thebibliography}{99}

    \bibitem{Moriondo2026VDT}
    L.~ Moriondo,
    \newblock {Vibrational Deduction Transformer (VDT) Leveraging Spectral methods, Theory of Sound and Associative Memory}.
    \newblock Zenodo.
    \newblock \url{https://doi.org/10.5281/zenodo.20816835}.

    \bibitem{Barontini2026}
    G.~Barontini,
    \newblock {Testing the problem of time with cold atoms},
    \newblock {\em Physical Review Research}, 8:L022047, 2026.
    \newblock doi: \href{https://doi.org/10.1103/1h9j-df4k}{10.1103/1h9j-df4k}.

    \bibitem{ConnesRovelli1994}
    A.~Connes and C.~Rovelli,
    \newblock {Von Neumann Algebra Automorphisms and Time-Thermodynamics Relation in General Covariant Quantum Theories},
    \newblock {\em Classical and Quantum Gravity}, 11:2899--2918, 1994.
    \newblock \href{https://arxiv.org/abs/gr-qc/9406019}{arXiv:gr-qc/9406019}.

    \bibitem{Moriondo2025JOSS}
    L.~Moriondo,
    \newblock {ArrowSpace: introducing Spectral Indexing for vector search},
    \newblock {\em Journal of Open Source Software}, 10(113):9002, 2025.
    \newblock doi: \href{https://doi.org/10.21105/joss.09002}{10.21105/joss.09002}.

    \bibitem{MoriondoAzizi2026}
    L.~Moriondo and I.~Azizi,
    \newblock {From Embedding Geometry to Spectral Search: Energy Dispersion Networks for Vector Retrieval},
    \newblock arXiv:2606.21535 [cs.IR], 2026.
    \newblock \href{https://arxiv.org/abs/2606.21535}{https://arxiv.org/abs/2606.21535}.

    \bibitem{Rombach2022}
    R.~Rombach, A.~Blattmann, D.~Lorenz, P.~Esser, and B.~Ommer,
    \newblock {High-Resolution Image Synthesis with Latent Diffusion Models},
    \newblock in {\em CVPR}, 2022.
    \newblock \href{https://arxiv.org/abs/2112.10752}{arXiv:2112.10752}.

    \bibitem{NicholDhariwal2021}
    A.~Nichol and P.~Dhariwal,
    \newblock {Improved Denoising Diffusion Probabilistic Models},
    \newblock in {\em ICML}, 2021.
    \newblock \href{https://arxiv.org/abs/2102.09672}{arXiv:2102.09672}.

    \bibitem{Moriondo2026ESDM_code}
    L.~Moriondo,
    \newblock {\textsc{\Large\faGithub} entropic-semantic-diffusion},
    \newblock \href{https://github.com/tuned-org-uk/entropic-semantic-diffusion}{Github}

    \bibitem{Moriondo2026ALDSC_code}
    L.~Moriondo,
    \newblock {\textsc{\Large\faGithub} arrowspace-latent-diffusion},
    \newblock \href{https://github.com/tuned-org-uk/arrowspace-latent-diffusion}{Github}

    \bibitem{Stancevic2025}
    D.~Stancevic, F.~Handke, and L.~Ambrogioni,
    \newblock {Entropic Time Schedulers for Generative Diffusion Models},
    \newblock arXiv:2504.13612 [cs.LG], 2025.
    \newblock \href{https://arxiv.org/abs/2504.13612}{arXiv:2504.13612}.

    \bibitem{Esteves2026}
    C.~Esteves and A.~Makadia,
    \newblock {Spectrally-Guided Diffusion Noise Schedules},
    \newblock in {\em ICML}, 2026.
    \newblock \href{https://arxiv.org/abs/2603.19222}{arXiv:2603.19222}.

    \end{thebibliography}

    \end{document}
